% قالب فارسی مقالهٔ حل تمرین؛ با XeLaTeX و XePersian بسازید.
\documentclass[12pt,a4paper]{article}
\usepackage{amsmath,amssymb,graphicx,booktabs,geometry,hyperref}
\geometry{margin=25mm}
\usepackage{xepersian}
% نام قلم را مطابق فونت نصب‌شدهٔ محیط تنظیم کنید.
\settextfont{Amiri}
\setlatintextfont{Latin Modern Roman}
\title{عنوان فارسی تمرین: حل و ارزیابی}
\author{نام پژوهشگر}
\date{تاریخ و نسخه}
\begin{document}
\maketitle
\begin{abstract}
هدف، روش، داده و مهم‌ترین نتیجهٔ واقعی را پس از اجرا در چند جمله گزارش کنید. نتیجهٔ فرضی ننویسید.
\end{abstract}
\noindent\textbf{واژگان کلیدی:} مسئله، روش عددی، اعتبارسنجی

\section{صورت مسئله}
منبع صورت سؤال، داده‌های معلوم، مجهولات، واحدها و دامنهٔ مسئله را بنویسید.
\section{فرض‌ها و فرمول‌بندی}
شرایط اولیه و مرزی و قرارداد نمادها را تعریف کنید. رابطهٔ زیر جای‌نگهدار است و باید حذف شود:
\begin{equation}
F(x)=0, \qquad x\in\Omega .
\label{eq:placeholder}
\end{equation}
\section{روش حل و بازتولید}
روش تحلیلی/عددی، پارامترها، نسخهٔ ابزار، معیار توقف و دستور اجرا را بنویسید.
\section{نتایج}
فقط پس از اجرای واقعی، جدول و شکل را با مقادیر واقعی وارد کنید. به هر کدام در متن ارجاع دهید.
% نمونهٔ ساختار جدول:
% \begin{table}[htbp]\centering
% \caption{عنوان دقیق با واحدهای ستون‌ها}\label{tab:results}
% \begin{tabular}{ll}\toprule کمیت & مقدار (واحد)\\\midrule ... & ...\\\bottomrule\end{tabular}
% \end{table}
% نمونهٔ شکل: \begin{figure}[htbp]\centering
% \includegraphics[width=.85\linewidth]{../results/figures/figure-01.pdf}
% \caption{عنوان، محورها، واحدها و شرایط تولید}\label{fig:result}\end{figure}
\section{بحث و اعتبارسنجی}
کنترل ابعادی، همگرایی/باقیمانده، مقایسه با مرجع مستقل در صورت وجود، خطا و محدودیت را تحلیل کنید.
\section{پاسخ نهایی}
به بندهای سؤال صریح و جداگانه پاسخ دهید.
\section{منابع}
منبع سؤال، روش و داده را با مشخصات قابل پیگیری ذکر کنید.
\appendix
\section{جزئیات تکمیلی}
در صورت نیاز داده یا مشتق‌گیری طولانی را بیاورید.
\end{document}
